\section{Conclusions}\label{sec:conclusions}

We have described a CLASS processing chain that uses only public data. It combines a simplified fundamental-parameter model, a CHIANTI solar spectrum at the GOES temperature of each exposure, the published filter transmission, a particle-background template and a Poisson likelihood fit with per-spectrum uncertainties and quality flags. Applied to \NRawFiles{} Level-1 exposures from \NMonths{} monthly archives, it gives \NMapped{} quality-controlled spectra and global 2\textdegree{} maps of Mg/Si and Al/Si with complete observed coverage. Our main findings are:
\begin{enumerate}
\item Element/Si ratios depend appreciably on the incident solar spectrum. Solar line emission carries most of the flux that excites Mg, Al and Si, and an isothermal spectrum at the GOES temperature brings the retrieved Mg/Si into the lunar range.
\item CLASS nightside spectra contain a solar-driven instrumental Al K$\alpha$ component when the Sun illuminates the spacecraft from behind. How it is treated in the dayside background model affects Al/Si.
\item With a GOES-based solar spectrum, Mg/Si and Al/Si can be recovered, whereas Ca/Si, Ti/Si and Fe/Si are dominated by solar-spectrum systematics.
\item Al/Si agrees with Lunar Prospector GRS ($r=\GrsAlR$) and with nine landing-site soils ($r=\SiteAlTwoR$), and several maria appear as distinct units. Mg/Si agrees with GRS ($r=\GrsMgR$); its landing-site comparison is inconclusive, as it is for GRS itself.
\item For three months, replacing the GOES-based spectrum by XSM-derived spectra changes the median Al/Si by a few per cent and Mg/Si by \XsmTwoMg{} to \XsmOneMg\%, depending on the assumed thermal structure. This supports Al/Si as the more robust product and identifies the solar spectrum as the main uncertainty for Mg/Si.
\end{enumerate}
The code, mapping-set results and maps are public. Extending the XSM-based treatment to all months, with free coronal abundances, is the most direct way to reduce the remaining systematic uncertainty.
